\section{Phase-averaged control and the distinguished zero phase}
\label{sec:phase-masks}

The angular formulation in \cref{sec:energy-cross-spectrum} is at the
integer-translation level.  There is a complementary factor-coordinate
test: modulate the source and quotient variables independently.  Define
\begin{align}
 \cH^\sharp(\alpha,\beta)
 ={}&\sum_{\substack{\ell>M\\k>1}}
 \Lambda(\ell)r_R(\ell k+h)W(\ell k/X)
 \e(\alpha\ell+\beta k).
 \label{eq:two-phase-hard-packet}
\end{align}
The original packet is the distinguished point
$\cH^\sharp(0,0)$.

For the dyadic estimates below, define
\begin{align}
 \cH^\sharp_{L,K}(\alpha,\beta)
 ={}&\sum_{\substack{\ell>M\\k>1}}
 \Lambda(\ell)r_R(\ell k+h)W(\ell k/X)
 \psi(\ell/L)\psi(k/K)
 \e(\alpha\ell+\beta k).
 \label{eq:dyadic-two-phase-packet}
\end{align}

We use the elementary divisor bound
\begin{equation}
 |r_R(n)|\ll_{\eps,h,W,\delta}X^\eps
 \qquad(n-h\asymp_WX).
 \label{eq:residual-pointwise-divisor-bound}
\end{equation}
Indeed, \eqref{eq:model-divisor-expansion-short}, the pointwise estimate
$|\lambda'_R(d)|\ll(d/\varphi(d))\log(2R/d)$, and the divisor bound give
$|\Lambda_R(n)|\ll_\eps X^\eps$, while
$\Lambda(n)\le\log n$.

\begin{theorem}[Factor-phase mean squares]
\label{thm:factor-phase-mean-squares}
For every $\eps>0$,
\begin{align}
 \int_\T|\cH^\sharp(\alpha,0)|^2\,d\alpha
 &\ll_{\eps,h,W,\delta}\frac{X^{2+\eps}}M,
 \label{eq:one-phase-mean-square}\\
 \int_\T\int_\T|\cH^\sharp(\alpha,\beta)|^2
 \,d\alpha\,d\beta
 &\ll_{\eps,h,W,\delta}X^{1+\eps}.
 \label{eq:two-phase-mean-square}
\end{align}
For a dyadic block with $LK\asymp_WX$,
\begin{align}
 \int_\T|\cH^\sharp_{L,K}(\alpha,0)|^2\,d\alpha
 &\ll_\eps\frac{X^{2+\eps}}L,
 \label{eq:dyadic-row-phase}\\
 \int_\T|\cH^\sharp_{L,K}(0,\beta)|^2\,d\beta
 &\ll_\eps\frac{X^{2+\eps}}K,
 \label{eq:dyadic-quotient-phase}\\
 \int_\T\int_\T|\cH^\sharp_{L,K}(\alpha,\beta)|^2
 \,d\alpha\,d\beta
 &\ll_\eps X^{1+\eps}.
 \label{eq:dyadic-two-phase}
\end{align}
\end{theorem}

\begin{proof}
Parseval in $\alpha$ gives
\begin{equation}
 \int_\T|\cH^\sharp(\alpha,0)|^2d\alpha
 =\sum_{\ell>M}\Lambda(\ell)^2
 \left|\sum_{k>1}r_R(\ell k+h)W(\ell k/X)\right|^2.
 \label{eq:one-phase-Parseval}
\end{equation}
For each $\ell$, the inner sum has $O_W(X/\ell)$ terms.  Apply
\eqref{eq:residual-pointwise-divisor-bound}, square, and use
\[
 \sum_{\ell>M}\frac{\Lambda(\ell)^2}{\ell^2}
 \ll\frac{\log(2M)}M.
\]
The logarithm is absorbed by $X^\eps$, proving
\eqref{eq:one-phase-mean-square}.

Double Parseval gives the exact diagonal sum
\begin{align*}
 \int_\T\int_\T|\cH^\sharp(\alpha,\beta)|^2
 \,d\alpha\,d\beta
 =\sum_{\ell>M,k>1}
 \Lambda(\ell)^2|r_R(\ell k+h)|^2|W(\ell k/X)|^2.
\end{align*}
The pointwise residual bound and
$\sum_{\ell\le y}\Lambda(\ell)^2\ll y\log(2y)$ make this
$O_\eps(X^{1+\eps})$.  Restricting the same arguments to
$\ell\asymp L,k\asymp K$ proves
\cref{eq:dyadic-row-phase,eq:dyadic-two-phase}.  For
\eqref{eq:dyadic-quotient-phase}, Parseval in $\beta$ and Cauchy--Schwarz
in $\ell$ give
\[
 \sum_{k\asymp K}
 \left|\sum_{\ell\asymp L}\Lambda(\ell)
 r_R(\ell k+h)W(\ell k/X)\psi(\ell/L)\right|^2
 \ll_\eps KL^2X^\eps.
\]
Here we used
$\sum_{\ell\asymp L}\Lambda(\ell)^2\ll L\log(2L)$ and absorbed the
logarithm into $X^\eps$.  Since $LK\asymp X$, this is
$O_\eps(X^{2+\eps}/K)$.
\end{proof}

Since $M=X^{1/2-o(1)}$, the root mean square in
\eqref{eq:one-phase-mean-square} is $X^{3/4+o(1)}$, and the double-phase
root mean square is $X^{1/2+o(1)}$.  More quantitatively, for every fixed
$\eta>0$, Chebyshev's inequality shows at each $X$ that the set of row
phases with $|\cH^\sharp(\alpha,0)|>X^{3/4+\eta}$ has measure
$O_\eta(X^{-2\eta+o(1)})$; the analogous double-phase exceptional measure
at threshold $X^{1/2+\eta}$ obeys the same bound.  Along a sufficiently
sparse dyadic subsequence, Borel--Cantelli gives the corresponding
almost-everywhere statements.  For independent Rademacher signs, the
expectations of the squared magnitudes equal the matching diagonal sums
in the Parseval calculations.

\begin{remark}[Why the theorem does not evaluate the gate]
The point $(\alpha,\beta)=(0,0)$ has Haar measure zero.  An $L^2$ estimate
over phases can coexist with a spike at that point.  In the original
arithmetic problem all rows are coherently unmodulated, so no phase
average is available.  The theorem identifies a phase-averaged mask family
with power cancellation; it does not transfer that cancellation to the
fixed all-positive direction.
\end{remark}

\subsection{Operator norms and actual arithmetic directions}

For a dyadic block define the residual matrix
\begin{equation}
 \mathsf R_{L,K}^{(h,R)}(\ell,k)=
 r_R(\ell k+h)W(\ell k/X)\psi(\ell/L)\psi(k/K).
 \label{eq:dyadic-residual-matrix}
\end{equation}
With $u_\ell=\Lambda(\ell)$ and, in the original TPC-15 coordinates,
$v_k=\beta_V(k)$, the block is $u^*\mathsf Rv$.  Standard divisor moments
give
\begin{equation}
 \|u\|_2^2\ll L\log(2L),
 \qquad
 \|v\|_2^2\ll K(\log(2K))^3.
 \label{eq:arithmetic-vector-norms}
\end{equation}
Thus a uniform operator estimate
\begin{equation}
 \|\mathsf R_{L,K}^{(h,R)}\|_{2\to2}
 =o\left(\frac{\sqrt X}{(\log X)^2}\right)
 \label{eq:sufficient-operator-norm}
\end{equation}
would make that block $o(X)$.

This is a sufficient condition for the actual arithmetic vectors.  It is
necessary and sufficient only for the full pair of $\ell^2$ mask balls,
because
\begin{equation}
 \sup_{\|x\|_2\le A,\ \|y\|_2\le B}
 |x^*\mathsf Ry|=AB\|\mathsf R\|_{2\to2}.
 \label{eq:operator-duality}
\end{equation}
The prime-pair problem asks for one specified coefficient
$u^*\mathsf Rv$, not the largest singular direction.  Replacing the
directional problem by \eqref{eq:sufficient-operator-norm} would therefore
ask for a stronger theorem than is logically necessary.

The available coarse Hilbert--Schmidt ledger is also not enough.  From
\eqref{eq:residual-pointwise-divisor-bound} one obtains the coarse bound
\begin{equation}
 \|\mathsf R_{L,K}^{(h,R)}\|_{\rm HS}^2
 \ll_\eps X^{1+\eps}.
 \label{eq:Hilbert-Schmidt-ledger}
\end{equation}
This bound controls total matrix energy but gives no small angle with the
vectors in \eqref{eq:arithmetic-vector-norms}.  The distinction mirrors
\cref{thm:residual-energy-law,cor:angular-corridor}.

\subsection{Boundary statement}

The results can now be placed in a precise order:
\begin{equation}
 \begin{array}{c}
 \text{finite local model and exact row drift}\\
 \Downarrow\\
 \text{all source rows below }X^{1/2-o(1)}\text{ are closed}\\
 \Downarrow\\
 \text{one factorable }(D,L,K)\text{ wedge is closed}\\
 \Downarrow\\
 \text{residual energy is large, but generic phases cancel}\\
 \Downarrow\\
 \text{the zero-phase fixed-shift directional coefficient remains.}
 \end{array}
 \label{eq:boundary-flow}
\end{equation}
The final arrow is not supplied by any theorem in the paper.  Existing
beyond-square-root distribution results require additional factorization
or coefficient structure; they do not by themselves overcome sieve
parity or select $h=2$ \citep{Maynard2025I,Polymath2014,FordMaynard2024}.
